\lecture{4.4}{IPTW with Discrete Data}
\label{lec:4-4}

\subsection{Estimating Propensity Scores by Hand}

The inverse probability weighted estimator is motivated by the identification result
\[
\E\{\Ya{a}\}
=
\E\left\{
\frac{\ind(A=a)}{g_a(W)}Y
\right\},
\qquad
g_a(W)=\Prob(A=a\given W).
\]
With discrete $W$, the propensity score can be estimated within each covariate
stratum. For binary treatment,
\[
g_{n,1}(w)
=
\frac{\sum_{i=1}^n \ind(A_i=1,W_i=w)}
       {\sum_{i=1}^n \ind(W_i=w)},
\]
and
\[
g_{n,0}(w)=1-g_{n,1}(w).
\]
Thus $g_{n,1}(w)$ is the sample proportion treated among individuals with $W=w$.

If $W$ is binary, this gives two treatment probabilities:
\[
g_{n,1}(0)=\frac{\#\{i:A_i=1,W_i=0\}}{\#\{i:W_i=0\}},
\qquad
g_{n,1}(1)=\frac{\#\{i:A_i=1,W_i=1\}}{\#\{i:W_i=1\}}.
\]

\subsection{Estimating Propensity Scores Using Regression}

The same quantities can be obtained from a saturated logistic regression for
treatment. With binary $W$,
\[
\Prob(A=1\given W=w)=\expit(\gamma_0+\gamma_1w).
\]
There are two parameters and two covariate strata, so the fitted probabilities can
match the empirical treatment proportions in the two strata.

\begin{implementationbox}
Saturated propensity score model for binary $W$:
\begin{lstlisting}[language=R]
fit_g <- glm(A ~ W, data = dat, family = binomial())

newW <- data.frame(W = c(0, 1))
g1_by_W <- predict(fit_g, newdata = newW, type = "response")
g0_by_W <- 1 - g1_by_W
\end{lstlisting}
\end{implementationbox}

For more complex discrete $W$, a saturated treatment model requires enough terms to
represent the treatment probability in every covariate stratum. When $W$ is large or
continuous, analysts usually use a lower-dimensional regression or machine learning
method to smooth across similar covariate values.

\subsection{Estimating the Joint Distribution}

The fully discrete IPTW plug-in estimator can be written by estimating the joint
distribution of $(W,A,Y)$. If $W$, $A$, and $Y$ are discrete, the empirical joint
probability is
\[
P_n(W=w,A=a,Y=y)
=
\frac{1}{n}\sum_{i=1}^n \ind(W_i=w,A_i=a,Y_i=y).
\]
Plugging this empirical distribution and the estimated propensity score into the IPTW
identification formula gives
\[
\psi_n(a)
=
\sum_{w,a',y}
\frac{\ind(a'=a)}{g_{n,a}(w)}
y\,P_n(W=w,A=a',Y=y).
\]
The summation uses $a'$ as a dummy treatment value to avoid confusion with the target
intervention level $a$.

\subsection{Computing Inverse Probability Weights}

For target treatment level $a$, the inverse probability weight for observation $i$ is
\[
\omega_i(a)=\frac{\ind(A_i=a)}{g_{n,a}(W_i)}.
\]
If an individual actually received treatment level $a$, their weight is the inverse of
the estimated probability of receiving that treatment. If they did not receive
treatment level $a$, their weight is zero for estimating $\E\{\Ya{a}\}$.

\begin{figure}[H]
\centering
\begin{tikzpicture}[
  node distance=1.35cm,
  box/.style={draw, rounded corners=2pt, minimum width=3.2cm, minimum height=.8cm},
  arr/.style={-{Latex[length=2.5mm]}, thick}
]
\node[box] (obs) {Observed treated\\ subgroup};
\node[box, right=of obs] (weight) {Weight by\\ $1/g_{n,1}(W)$};
\node[box, right=of weight] (pseudo) {Represents full\\ covariate mix};
\draw[arr] (obs) -- (weight);
\draw[arr] (weight) -- (pseudo);
\node[box, below=1.1cm of obs] (obs0) {Observed untreated\\ subgroup};
\node[box, right=of obs0] (weight0) {Weight by\\ $1/g_{n,0}(W)$};
\node[box, right=of weight0] (pseudo0) {Represents full\\ covariate mix};
\draw[arr] (obs0) -- (weight0);
\draw[arr] (weight0) -- (pseudo0);
\end{tikzpicture}
\caption{IPTW reweights each observed treatment group to represent the target population.}
\label{fig:ch4-iptw-reweighting}
\end{figure}

Extreme weights arise when $g_{n,a}(W_i)$ is close to zero. Such observations can
dominate the estimate, which is why positivity is not merely a technical detail.

\subsection{Estimating $\E\{\Yone\}$, $\E\{\Yzero\}$, and the ATE}

Using the empirical joint distribution, the estimator simplifies to the individual-level
sample average
\[
\psi_n(a)
=
\frac{1}{n}\sum_{i=1}^n
\frac{\ind(A_i=a)}{g_{n,a}(W_i)}Y_i.
\]
The derivation follows by substituting the empirical distribution:
\[
\begin{aligned}
\sum_{w,a',y}
&\frac{\ind(a'=a)}{g_{n,a}(w)}
y\,P_n(W=w,A=a',Y=y)\\
&=
\sum_{w,a',y}
\frac{\ind(a'=a)}{g_{n,a}(w)}
y\left\{
\frac{1}{n}\sum_{i=1}^n
\ind(W_i=w,A_i=a',Y_i=y)
\right\}\\
&=
\frac{1}{n}\sum_{i=1}^n
\sum_{w,a',y}
\frac{\ind(a'=a)}{g_{n,a}(w)}
y\ind(W_i=w,A_i=a',Y_i=y)\\
&=
\frac{1}{n}\sum_{i=1}^n
\frac{\ind(A_i=a)}{g_{n,a}(W_i)}Y_i.
\end{aligned}
\]
Therefore,
\[
\psi_n(1)
=
\frac{1}{n}\sum_{i=1}^n
\frac{\ind(A_i=1)}{g_{n,1}(W_i)}Y_i,
\]
\[
\psi_n(0)
=
\frac{1}{n}\sum_{i=1}^n
\frac{\ind(A_i=0)}{g_{n,0}(W_i)}Y_i,
\]
and
\[
\widehat{\ATE}_{\IPTW}=\psi_n(1)-\psi_n(0).
\]

Using the same binary-$W$ summaries from empirical g-computation, suppose
\[
g_{n,1}(0)=0.476,\qquad g_{n,0}(0)=0.524,
\]
and
\[
g_{n,1}(1)=0.486,\qquad g_{n,0}(1)=0.514.
\]
The IPTW calculation weights each observed outcome by the inverse probability of
the treatment actually received. When the propensity scores are empirical stratum
proportions, the weighted treated mean reproduces the same standardized quantity
as empirical g-computation:
\[
\psi_n(1)=2.530,\qquad \psi_n(0)=2.392,
\]
and therefore
\[
\widehat{\ATE}_{\IPTW}=2.530-2.392=0.138.
\]
This numerical equality is a special property of the fully empirical discrete
estimator. It should not be expected once lower-dimensional regression smoothing is
used for $\Qbar_n$ or $g_n$.

\begin{implementationbox}
Regression-based IPTW in R:
\begin{lstlisting}[language=R]
fit_g <- glm(A ~ W, data = dat, family = binomial())

g1 <- predict(fit_g, type = "response")
g0 <- 1 - g1

psi1 <- mean((dat$A == 1) / g1 * dat$Y)
psi0 <- mean((dat$A == 0) / g0 * dat$Y)
ate  <- psi1 - psi0
\end{lstlisting}
\end{implementationbox}

\begin{estimationbox}
By-hand IPTW with discrete data:
\begin{enumerate}
\item Estimate $g_{n,1}(w)$ and $g_{n,0}(w)$ in every covariate stratum.
\item For each observation, compute $\omega_i(1)=\ind(A_i=1)/g_{n,1}(W_i)$ and
$\omega_i(0)=\ind(A_i=0)/g_{n,0}(W_i)$.
\item Compute $\psi_n(1)=n^{-1}\sum_i \omega_i(1)Y_i$.
\item Compute $\psi_n(0)=n^{-1}\sum_i \omega_i(0)Y_i$.
\item Subtract to estimate the ATE.
\end{enumerate}
\end{estimationbox}
